\chapter{Scaling Up the Architecture}
\label{ch:scaling}

BabyGPT has 6 layers, 6 heads, and a 384-dimensional embedding --- about 10.7 million parameters.  GPT-2, the model described in \paper{Language Models are Unsupervised Multitask Learners} (Radford et al., 2019), scales every one of these axes.  In this chapter we walk through exactly what changes, derive the parameter count, and build the canonical GPT-2 model in PyTorch.

The architecture is structurally identical to BabyGPT --- pre-norm transformer blocks with causal self-attention --- just bigger.  If you understood the basic tutorial, you already understand GPT-2.  The differences are in the numbers.


\section{GPT-2 Model Sizes}

The GPT-2 paper defines four model sizes:

\begin{center}
\begin{tabular}{lrrrrr}
\toprule
Model & Layers & Heads & $d_{\text{model}}$ & Context & Parameters \\
\midrule
GPT-2 Small  & 12 & 12 & 768   & 1\,024 & 117M \\
GPT-2 Medium & 24 & 16 & 1\,024 & 1\,024 & 345M \\
GPT-2 Large  & 36 & 20 & 1\,280 & 1\,024 & 762M \\
GPT-2 XL     & 48 & 25 & 1\,600 & 1\,024 & 1\,542M \\
\bottomrule
\end{tabular}
\end{center}

All four share the same architecture --- they differ only in these hyperparameters.  Notice that the head dimension ($d_{\text{model}} / n_{\text{heads}}$) is always 64, matching BabyGPT.

We target the smallest variant.  It is large enough to produce coherent text but small enough to train on a modest GPU cluster.

\tip{The parameter counts in the table are from the original paper, which used a 50\,257-token vocabulary and weight tying (sharing the embedding and output matrices).  Our implementation uses a 32\,768-token vocabulary with no weight tying, so our count will differ slightly.  We will derive the exact number below.}


\section{What Changes from BabyGPT}

\begin{center}
\begin{tabular}{lrr}
\toprule
Hyperparameter & BabyGPT & MicroGPT \\
\midrule
\code{n\_layer} & 6 & 12 \\
\code{n\_head}  & 6 & 12 \\
\code{n\_embd}  & 384 & 768 \\
\code{context / sequence\_len} & 256 & 1\,024 \\
\code{vocab\_size} & 65 & 32\,768 \\
\code{dropout}  & 0.2 & 0.0 \\
Weight tying    & Yes & No \\
\bottomrule
\end{tabular}
\end{center}

\textbf{No dropout.}  BabyGPT uses 0.2 dropout because TinyShakespeare is tiny (1\,MB) and overfitting is a real concern.  At GPT-2 scale the training set is billions of tokens --- the model will never see the same data twice, so dropout is unnecessary.

\textbf{No weight tying.}  BabyGPT shares the token embedding matrix with the output head (\code{lm\_head}).  The GPT-2 paper does not tie these weights.  With separate matrices, the embedding can specialize in turning token IDs into vectors while the output head specializes in turning vectors back into logits.  The cost is an extra $\text{vocab\_size} \times d_{\text{model}}$ parameters.

\textbf{Vocab padding.}  GPU tensor cores operate on blocks of 8 or 16 elements.  If the vocabulary size is not a multiple of 64, the last block is partially empty and wastes compute.  We pad $32\text{,}768$ up to the nearest multiple of 64 --- which is already $32\text{,}768 = 2^{15}$, so no padding is needed in our case.  The code handles non-aligned sizes automatically.


\section{Parameter Count}

Let us derive the parameter count for MicroGPT ($d = 768$, $L = 12$, $h = 12$, $V = 32\text{,}768$, $T = 1\text{,}024$).

\textbf{Token embedding} \code{wte}: $V \times d = 32\text{,}768 \times 768 = 25\text{,}165\text{,}824$.

\textbf{Position embedding} \code{wpe}: $T \times d = 1\text{,}024 \times 768 = 786\text{,}432$.

\textbf{Per transformer block} (repeated $L = 12$ times):
\begin{itemize}
  \item \code{W\_Q}, \code{W\_K}, \code{W\_V}: each $d \times d = 589\text{,}824$.  Three projections: $3 \times 589\text{,}824 = 1\text{,}769\text{,}472$.
  \item \code{c\_proj} (attention output): $d \times d = 589\text{,}824$.
  \item \code{c\_fc} (FFN expand): $d \times 4d = 768 \times 3\text{,}072 = 2\text{,}359\text{,}296$.
  \item \code{c\_proj} (FFN contract): $4d \times d = 2\text{,}359\text{,}296$.
  \item \code{ln\_1}, \code{ln\_2}: each has $d = 768$ learnable weights (no bias).  Total: $1\text{,}536$.
  \item \textbf{Block total}: $1\text{,}769\text{,}472 + 589\text{,}824 + 2\times 2\text{,}359\text{,}296 + 1\text{,}536 = 7\text{,}079\text{,}424$.
\end{itemize}

\textbf{All 12 blocks}: $12 \times 7\text{,}079\text{,}424 = 84\text{,}953\text{,}088$.

\textbf{Final LayerNorm} \code{ln\_f}: $768$.

\textbf{Output head} \code{lm\_head}: $d \times V = 768 \times 32\text{,}768 = 25\text{,}165\text{,}824$.

\textbf{Grand total}: $25\text{,}165\text{,}824 + 786\text{,}432 + 84\text{,}953\text{,}088 + 768 + 25\text{,}165\text{,}824 = \mathbf{136\text{,}071\text{,}936}$ ($\approx 136$M).

\tip{This is higher than the commonly cited ``GPT-2 124M'' because (1) the original paper used weight tying, which saves $V \times d \approx 25$M parameters, and (2) the original vocabulary was 50\,257 with weight tying, while ours is 32\,768 without tying.  If we tied weights, our count would be $\sim$111M.  The transformer body (blocks + norms) accounts for 85M regardless --- that is the core model, and it is the same scale as the original.}


\section{Weight Initialization}

GPT-2 initializes all weights from $\mathcal{N}(0, 0.02)$ with one important exception: \textbf{residual projection} weights (\code{c\_proj} in both attention and FFN) are scaled by $1/\sqrt{2L}$ where $L$ is the number of layers.

Why?  Each layer adds a residual contribution.  After $L$ layers, the variance of the residual stream grows as $\sim L$ if each contribution has unit variance.  Scaling the initial variance of residual projections by $1/(2L)$ keeps the total variance bounded.  The factor of 2 accounts for the two residual additions per block (attention + FFN).

\begin{align*}
\text{Normal weights:} \quad & w \sim \mathcal{N}(0, 0.02) \\
\text{Residual projections:} \quad & w \sim \mathcal{N}\!\left(0, \frac{0.02}{\sqrt{2L}}\right)
\end{align*}


\section{Implementation}

The config is a simple dataclass:

\filetag{microgpt/config.py}
\begin{lstlisting}
"""MicroGPT model configuration."""

from dataclasses import dataclass


@dataclass
class GPTConfig:
    """GPT-2 124M configuration.

    This is the smallest GPT-2 variant from Radford et al. (2019).
    Compared to BabyGPT (6 layers, 6 heads, 384 dim, 256 context, 65 vocab):
    every axis is scaled up, and character-level tokenization is replaced by BPE.
    """
    sequence_len: int = 1024     # context window (was 256 in BabyGPT)
    vocab_size: int = 32768      # BPE vocabulary (was 65)
    n_layer: int = 12            # transformer blocks (was 6)
    n_head: int = 12             # attention heads (was 6)
    n_embd: int = 768            # embedding dimension (was 384)
    dropout: float = 0.0         # no dropout at scale (was 0.2)
    bias: bool = False           # no bias in Linear/LayerNorm (same as BabyGPT)
\end{lstlisting}

The model is structurally identical to BabyGPT --- only the numbers change, plus the removal of weight tying and the addition of vocab padding:

\filetag{microgpt/model.py}
\begin{lstlisting}
"""MicroGPT: a GPT-2 124M decoder-only transformer language model.

This is the canonical GPT-2 architecture from Radford et al. (2019):
pre-norm transformer blocks with LayerNorm, GELU activations, learned
positional embeddings, and no weight tying.  Chapter 3 of the tutorial
upgrades this to modern techniques (RMSNorm, RoPE, GQA, relu-squared).
"""

import math

import torch
import torch.nn as nn
import torch.nn.functional as F

from .config import GPTConfig


# ---------------------------------------------------------------------------
# Building blocks
# ---------------------------------------------------------------------------

class LayerNorm(nn.Module):
    """LayerNorm with optional bias (identical to BabyGPT).

    LayerNorm(x) = gamma * (x - mean) / sqrt(var + eps) + beta
    """

    def __init__(self, n_embd: int, bias: bool = False, eps: float = 1e-5) -> None:
        super().__init__()
        self.weight = nn.Parameter(torch.ones(n_embd))
        self.bias = nn.Parameter(torch.zeros(n_embd)) if bias else None
        self.eps = eps

    def forward(self, x: torch.Tensor) -> torch.Tensor:
        mean = x.mean(dim=-1, keepdim=True)
        var = x.var(dim=-1, keepdim=True, unbiased=False)
        x = (x - mean) / torch.sqrt(var + self.eps)
        x = self.weight * x
        if self.bias is not None:
            x = x + self.bias
        return x


class CausalSelfAttention(nn.Module):
    """Multi-head self-attention with causal mask.

    Identical structure to BabyGPT but scaled up: 12 heads of 64 dims each
    (768 / 12 = 64).  Uses learned positional embeddings passed through the
    residual stream --- positions are not handled inside attention.
    """

    def __init__(self, config: GPTConfig) -> None:
        super().__init__()
        assert config.n_embd % config.n_head == 0
        self.n_head = config.n_head
        self.n_embd = config.n_embd
        self.head_dim = config.n_embd // config.n_head

        self.W_Q = nn.Linear(config.n_embd, config.n_embd, bias=config.bias)
        self.W_K = nn.Linear(config.n_embd, config.n_embd, bias=config.bias)
        self.W_V = nn.Linear(config.n_embd, config.n_embd, bias=config.bias)
        self.c_proj = nn.Linear(config.n_embd, config.n_embd, bias=config.bias)
        self.attn_dropout = nn.Dropout(config.dropout)
        self.resid_dropout = nn.Dropout(config.dropout)

    def forward(self, x: torch.Tensor) -> torch.Tensor:
        B, T, C = x.shape

        q = self.W_Q(x).view(B, T, self.n_head, self.head_dim).transpose(1, 2)
        k = self.W_K(x).view(B, T, self.n_head, self.head_dim).transpose(1, 2)
        v = self.W_V(x).view(B, T, self.n_head, self.head_dim).transpose(1, 2)

        # Scaled dot-product attention with causal mask
        scores = torch.matmul(q, k.transpose(-2, -1)) / math.sqrt(self.head_dim)
        mask = torch.tril(torch.ones(T, T, device=x.device, dtype=torch.bool))
        scores = scores.masked_fill(~mask, float("-inf"))
        weights = F.softmax(scores, dim=-1)
        weights = self.attn_dropout(weights)
        out = torch.matmul(weights, v)

        # Reassemble heads
        out = out.transpose(1, 2).contiguous().view(B, T, C)
        return self.resid_dropout(self.c_proj(out))


class FeedForward(nn.Module):
    """Position-wise FFN with GELU activation.

    FFN(x) = W2 * GELU(W1 * x)
    Hidden dim is 4x embedding dim (expand then contract).
    """

    def __init__(self, config: GPTConfig) -> None:
        super().__init__()
        self.c_fc = nn.Linear(config.n_embd, 4 * config.n_embd, bias=config.bias)
        self.c_proj = nn.Linear(4 * config.n_embd, config.n_embd, bias=config.bias)
        self.dropout = nn.Dropout(config.dropout)

    def forward(self, x: torch.Tensor) -> torch.Tensor:
        x = self.c_fc(x)
        x = F.gelu(x)
        x = self.c_proj(x)
        return self.dropout(x)


class TransformerBlock(nn.Module):
    """Pre-norm transformer block (identical structure to BabyGPT).

    x = x + Attention(LayerNorm(x))
    x = x + FFN(LayerNorm(x))
    """

    def __init__(self, config: GPTConfig) -> None:
        super().__init__()
        self.ln_1 = LayerNorm(config.n_embd, bias=config.bias)
        self.attn = CausalSelfAttention(config)
        self.ln_2 = LayerNorm(config.n_embd, bias=config.bias)
        self.ffn = FeedForward(config)

    def forward(self, x: torch.Tensor) -> torch.Tensor:
        x = x + self.attn(self.ln_1(x))
        x = x + self.ffn(self.ln_2(x))
        return x


# ---------------------------------------------------------------------------
# Full model
# ---------------------------------------------------------------------------

class MicroGPT(nn.Module):
    """Decoder-only transformer language model (GPT-2 124M).

    Compared to BabyGPT:
    - Larger: 12 layers, 12 heads, 768 dim, 1024 context, 32K vocab
    - No weight tying: lm_head has its own parameters (follows GPT-2 paper)
    - Vocab padding: vocab_size is rounded up to a multiple of 64 for
      tensor-core alignment (the extra embeddings are never used)
    """

    def __init__(self, config: GPTConfig) -> None:
        super().__init__()
        self.config = config

        # Round vocab_size up to nearest multiple of 64 for tensor-core efficiency
        padded_vocab = ((config.vocab_size + 63) // 64) * 64

        self.transformer = nn.ModuleDict(dict(
            wte=nn.Embedding(padded_vocab, config.n_embd),
            wpe=nn.Embedding(config.sequence_len, config.n_embd),
            drop=nn.Dropout(config.dropout),
            blocks=nn.ModuleList(
                [TransformerBlock(config) for _ in range(config.n_layer)]
            ),
            ln_f=LayerNorm(config.n_embd, bias=config.bias),
        ))
        self.lm_head = nn.Linear(config.n_embd, padded_vocab, bias=False)

        # No weight tying (unlike BabyGPT).  GPT-2 paper uses separate
        # embedding and output matrices.

        # Initialize weights
        self.apply(self._init_weights)
        # Scaled init for residual projections (GPT-2 convention):
        # scale by 1/sqrt(2 * n_layer) so residual variance stays bounded.
        for name, p in self.named_parameters():
            if name.endswith("c_proj.weight"):
                nn.init.normal_(p, mean=0.0, std=0.02 / math.sqrt(2 * config.n_layer))

    def _init_weights(self, module: nn.Module) -> None:
        if isinstance(module, nn.Linear):
            nn.init.normal_(module.weight, mean=0.0, std=0.02)
            if module.bias is not None:
                nn.init.zeros_(module.bias)
        elif isinstance(module, nn.Embedding):
            nn.init.normal_(module.weight, mean=0.0, std=0.02)

    def forward(
        self,
        tokens: torch.Tensor,
        targets: torch.Tensor | None = None,
    ) -> tuple[torch.Tensor, torch.Tensor | None]:
        B, T = tokens.shape
        assert T <= self.config.sequence_len

        positions = torch.arange(T, device=tokens.device)
        x = self.transformer.wte(tokens) + self.transformer.wpe(positions)
        x = self.transformer.drop(x)

        for block in self.transformer.blocks:
            x = block(x)

        x = self.transformer.ln_f(x)
        logits = self.lm_head(x)  # (B, T, padded_vocab)
        # Crop logits to actual vocab_size (remove padding tokens)
        logits = logits[..., :self.config.vocab_size]

        loss = None
        if targets is not None:
            loss = F.cross_entropy(
                logits.view(-1, logits.size(-1)),
                targets.view(-1),
            )
        return logits, loss

    @torch.no_grad()
    def generate(
        self,
        tokens: torch.Tensor,
        max_new_tokens: int,
        temperature: float = 1.0,
        top_k: int | None = None,
    ) -> torch.Tensor:
        """Autoregressive generation (identical to BabyGPT)."""
        for _ in range(max_new_tokens):
            tokens_cropped = (
                tokens if tokens.size(1) <= self.config.sequence_len
                else tokens[:, -self.config.sequence_len:]
            )
            logits, _ = self(tokens_cropped)
            logits = logits[:, -1, :] / temperature

            if top_k is not None:
                topk_vals, _ = torch.topk(logits, min(top_k, logits.size(-1)))
                logits[logits < topk_vals[:, [-1]]] = float("-inf")

            probs = F.softmax(logits, dim=-1)
            next_token = torch.multinomial(probs, num_samples=1)
            tokens = torch.cat([tokens, next_token], dim=1)

        return tokens

    def count_parameters(self) -> int:
        """Trainable parameters (all are unique --- no weight tying)."""
        return sum(p.numel() for p in self.parameters())
\end{lstlisting}

\checkpoint{Instantiate the model and verify the parameter count:}

\shelltag
\begin{lstlisting}[language=bash]
uv run python -c "
from microgpt.config import GPTConfig
from microgpt.model import MicroGPT
model = MicroGPT(GPTConfig())
print(f'Parameters: {model.count_parameters():,}')
"
\end{lstlisting}

\outputtag
\begin{lstlisting}[language={}]
Parameters: 136,071,936
\end{lstlisting}


\section{Comparing with BabyGPT}

Here is how parameter mass is distributed:

\begin{center}
\begin{tabular}{lrr}
\toprule
Component & BabyGPT & MicroGPT \\
\midrule
Token embedding (\code{wte})     & 24\,960   & 25\,165\,824 \\
Position embedding (\code{wpe})  & 98\,304   & 786\,432 \\
Transformer blocks               & 10\,621\,440 & 84\,953\,088 \\
Final LayerNorm                   & 384       & 768 \\
Output head (\code{lm\_head})    & 0 (tied)  & 25\,165\,824 \\
\midrule
\textbf{Total}                    & \textbf{10.7M} & \textbf{136.1M} \\
\bottomrule
\end{tabular}
\end{center}

At BabyGPT scale the embedding is negligible (0.2\% of total parameters).  At GPT-2 scale, the two embedding matrices (token + output) account for 37\% of parameters.  The transformer blocks --- the part that actually ``thinks'' --- account for 62\%.

\tip{This is why scaling laws focus on non-embedding parameters.  Doubling the vocabulary doubles the embedding cost but does not make the model smarter.  Doubling the number of layers or the embedding dimension adds capacity where it matters.}


\section{Summary}

We built the canonical GPT-2 124M model by scaling up BabyGPT:

\begin{enumerate}
  \item \textbf{Bigger dimensions}: 12 layers, 12 heads, 768 embedding dim, 1\,024 context.
  \item \textbf{BPE vocabulary}: 32\,768 tokens instead of 65 characters.
  \item \textbf{No weight tying}: separate embedding and output head.
  \item \textbf{No dropout}: unnecessary with large training sets.
  \item \textbf{Vocab padding}: round up for tensor-core alignment.
\end{enumerate}

The architecture is identical to BabyGPT --- the same pre-norm transformer block with LayerNorm, GELU, and learned positional embeddings.  In the next chapter we replace several of these components with modern alternatives that improve training stability, inference efficiency, and model quality.
